Continuous integration and continuous deployment (CI/CD) pipelines are now part of everyday software
development. In modern repositories, every pull request, dependency update, and workflow change can trigger a
GitHub Actions run that either confirms the health of the system or exposes a failure that blocks development.
Although these failures are valuable signals, they are also expensive in practice. Developers must inspect raw
logs, determine whether the failure is caused by configuration drift, dependency mismatch, test regressions,
environment problems, or transient infrastructure issues, and then decide whether the right response is to
patch code, re-run the workflow, or escalate the issue for human review. This triage work is repetitive, time
consuming, and often urgent because failed CI blocks merges, delays releases, and interrupts developer flow.

Recent progress in large language models (LLMs) has made it possible to automate parts of this process
\cite{logsage2025,cirepairbench2026}.
Log-analysis systems can summarize failures and suggest likely causes, while repository-aware repair systems can
attempt patches and validate them by re-running CI. However, the practical problem is not simply whether a
model can generate a patch. In real engineering workflows, an inappropriate automatic fix can be worse than no
fix at all. A tool that edits the wrong file, suppresses a failing test, disables a check, or proposes a risky
change for a transient failure may turn a local CI incident into a broader reliability problem. For this
reason, the central question of this thesis is not ``can an LLM make CI green?'' but rather ``what is the right
action for a given CI failure, and how can that action be taken safely?''

This thesis presents \textbf{PatchOps}, a repository-aware AI agent for GitHub Actions failures that combines
failure diagnosis, repository context retrieval, patch generation, local CI re-validation, and explicit
safety-gated decision making. Instead of treating every failure as an invitation to generate a patch, PatchOps
frames CI repair as a decision under uncertainty. Given a failed workflow run and its repository context, the
system decides among several actions: open a patch pull request, present a preview for human review, recommend
a re-run for likely transient failures, escalate risky cases, or abstain when evidence is insufficient. This
decision-centric framing is the thesis backbone.

\section{Research Context and Gap}

The research landscape already contains strong prior work on CI failure analysis and repository-level repair.
In particular, CI-Repair-Bench establishes CI-validated repository repair as a benchmarked task
\cite{cirepairbench2026}, while industrial systems such as LogSage demonstrate that LLM-based diagnosis can be
useful at scale \cite{logsage2025}. As a result, PatchOps does not claim to be the first system to perform
CI-verified repair, nor does it claim to surpass the state of the art in raw repair success. Such claims would be
inaccurate and would obscure the actual research opportunity.

The genuine gap addressed by this thesis is different. Existing CI repair systems largely optimize for whether a
patch can be generated and whether CI turns green after that patch is applied. They do not place equal emphasis
on whether the agent should patch in the first place, how often it should refuse to act, whether a green result
is obtained honestly, or whether the selected patch is even targeted at the right part of the repository. In
other words, the literature has established CI repair as a task, but it has not adequately framed repair as a
\emph{risk-sensitive decision process}. This is where PatchOps contributes: it adds a safety and abstention
layer on top of an already established repair setting, and evaluates that layer explicitly rather than treating
it as an implementation detail.

\section{PatchOps Thesis Position}

PatchOps demonstrably repairs CI failures end to end: it produces CI-verified fixes that turn failed workflows
green, shown on real GitHub pull requests whose checks passed after the patch. The repair capability is therefore
not in question. The position taken in this thesis builds on that capability rather than disputing it. PatchOps is
claimed as a reproducible, safety-gated, abstention-aware CI-repair agent that is dependable when failures are
log-localizable and cautious when they are not. The main contribution is consequently not raw patch volume,
aggressive automation, or a higher repair rate than prior work; it is the engineering of \emph{trustworthy}
repair --- knowing \emph{when not to act} --- together with a closed validation loop that checks proposed fixes
through CI-style re-execution before they are treated as credible.

The current evidence base in the project supports four findings that define the thesis:

\begin{enumerate}
\item Fault localization is the binding constraint, and within the methods tested it is not closed by retrieval
alone. Across four strategies evaluated in this work --- lexical BM25, LLM reranking, semantic embedding
retrieval, and bounded agentic (ReAct) exploration --- recall improves only modestly and plateaus below the
target bar needed for robust repair. We therefore argue, for our setting, that simple retrieval is insufficient
and that the residual difficulty appears to be one of repository reasoning; we do not claim this as a universal
property of CI repair.
\item Better localization is necessary but not sufficient for successful repair. Full-repository retrieval can
remove patch-application bottlenecks without producing corresponding gains in end-to-end Pass@1 repair success.
\item A green CI result is not automatically a correct result. Some patches can be off-target or reward-hacking
in nature, so honesty and plausibility checks are required in addition to validation.
\item Abstention cannot be left to the base agent alone. A tool-using agent may rarely self-decline even when it
should, so abstention must be engineered through explicit selective controls, safety gates, and the learned
calibrated act-versus-hold gate introduced in the final system.
\end{enumerate}

These findings are valuable even when they are negative or partially negative. A central strength of the
PatchOps study is that it does not treat honest negative results as failures to hide. Instead, they are used to
separate what the system already does well from what remains unsolved in CI repair, especially semantic repair
after the correct file has already been identified.

\section{Problem Motivation}

The motivation for PatchOps comes from the mismatch between the operational reality of CI failures and the way
automated repair is often discussed. Not all failures should be patched. Some runs fail because of flaky tests,
network issues, missing credentials, permission errors, deployment conditions, or ambiguous semantic
regressions. In such cases, blindly generating a patch is either unhelpful or unsafe. A production-relevant
repair assistant therefore needs at least three properties.

First, it must be \textbf{repository-aware}. CI failures are rarely understandable from logs alone. Workflow
files, dependency manifests, configuration files, and the surrounding codebase provide crucial context for
understanding what actually broke.

Second, it must be \textbf{validation-aware}. A generated patch should not be treated as useful merely because
it looks plausible in text. Practical repair demands a feedback loop in which the patch is applied and the
relevant CI commands are re-executed.

Third, it must be \textbf{safety-aware}. The system must recognize risky conditions, avoid touching forbidden or
sensitive files, and prefer preview, re-run, or escalation modes over automation when evidence is weak or the
stakes are high.

PatchOps is designed around these three requirements. The resulting system is not a general autonomous software
engineer; it is a scoped CI-repair agent whose primary value lies in disciplined decision making under noisy and
incomplete evidence.

\section{Research Questions}

This thesis is organized around five fixed research questions:

\begin{enumerate}
\item Does repository-aware context improve CI/CD failure diagnosis and fault localization compared with log-only
prompting?
\item Does PatchOps generate patches that apply cleanly and improve CI repair success compared with simpler LLM
baselines?
\item Does a safety-gated decision policy reduce unsafe or inappropriate repair attempts?
\item Can PatchOps distinguish repairable failures from transient or flaky failures where re-running is a better
action than patching?
\item What is the cost and latency of running PatchOps in a realistic GitHub workflow?
\end{enumerate}

These questions keep the work grounded. They separate diagnosis from repair, repair from safety, and automation
from practicality. They also ensure that the thesis does not collapse into a vague ``AI for DevOps'' claim.

\section{Contributions}

The contributions of this thesis are summarized as follows:

\begin{enumerate}
\item \textbf{A system contribution:} PatchOps, an end-to-end repository-aware repair agent for GitHub Actions
failures that integrates diagnosis, localization, patch planning, generation, validation, and decision making.
\item \textbf{A methodological contribution:} a staged repair architecture that treats CI repair as a
risk-sensitive decision process rather than a patch-generation task alone.
\item \textbf{A safety contribution:} explicit support for patch, preview, re-run, escalate, and abstain actions,
along with safety-gated controls over automated behavior.
\item \textbf{An evaluation contribution:} a benchmarked study that reports not only diagnosis and repair
metrics, but also unsafe-edit rate, abstention accuracy, calibration-oriented selective metrics, and
patch-plausibility indicators.
\item \textbf{An empirical contribution:} a set of findings showing that localization improvements alone do not
solve CI repair, that green CI outcomes require honesty checks, and that abstention must be engineered rather
than assumed.
\end{enumerate}

\section{Thesis Organization}

The remainder of this thesis is organized as follows. Chapter 2 reviews the background and related work on LLM
agents, CI failure diagnosis, automated program repair, safety, calibration, and abstention. Chapter 3 defines
the problem, scope, and safety requirements of PatchOps. Chapter 4 presents the system design and architecture,
including the decision pipeline and validation loop. Chapter 5 describes the implementation of the PatchOps
engine and its supporting modules. Chapter 6 explains the experimental methodology, datasets, baselines, and
evaluation metrics. Chapter 7 reports the empirical results, structured around the thesis's four headline
findings. Chapter 8 discusses the broader meaning of these findings and their implications for practical CI
repair agents. Chapter 9 examines threats to validity and current limitations. Chapter 10 concludes the thesis
and outlines future research directions, with particular emphasis on semantic repair and stronger verification.

\section{Chapter Conclusion}

This chapter has introduced the PatchOps thesis, its motivation, research questions, and contributions. It has
also established the central framing used throughout the document: CI repair is not only a patch-generation
problem, but a risk-sensitive decision problem requiring repository context, validation, safety gates, and
engineered abstention. The next chapter situates this framing within the related literature.
